% Auto-generated from raw.txt
% Usage:
%   \vocabentry{word}{/ipa/ (pos)}{definition}{example}{note}

\vocabentry{Abduction}
{/æbˈdʌkʃn/ (n)}
{The action of forcibly taking someone away against their will.}
{The police are investigating the abduction of a young child from the local park.}
{Đây là một danh từ dùng để chỉ vụ bắt cóc, thường dùng trong ngữ cảnh pháp lý.}

\vocabentry{Acquit}
{/əˈkwɪt/ (v)}
{Free someone from a criminal charge by a verdict of not guilty.}
{The jury decided to acquit the defendant due to a lack of concrete evidence.}
{Đây là một động từ dùng để chỉ việc tuyên trắng án cho một ai đó.}

\vocabentry{Arson}
{/ˈɑːrsn/ (n)}
{The criminal act of deliberately setting fire to property.}
{Several buildings were destroyed in what the fire department believes was an act of arson.}
{Đây là một danh từ dùng để chỉ tội cố ý gây hỏa hoạn/đốt nhà.}

\vocabentry{Assault}
{/əˈsɔːlt/ (n/v)}
{A physical attack; to make a physical attack on someone.}
{The victim was hospitalized after a brutal assault in the city center last night.}
{Đây là một danh từ/động từ dùng để chỉ hành vi hành hung hoặc tấn công người khác.}

\vocabentry{Burglary}
{/ˈbɜːrɡləri/ (n)}
{Entry into a building illegally with intent to commit a crime, especially theft.}
{The family returned from vacation to find that their house had been a victim of burglary.}
{Đây là một danh từ dùng để chỉ tội trộm cướp bằng cách đột nhập vào nhà.}

\vocabentry{Capital punishment}
{/ˌkæpɪtl ˈpʌnɪʃmənt/ (n)}
{The legally authorized killing of someone as punishment for a crime.}
{Many countries have abolished capital punishment, arguing that it is inhumane.}
{Đây là một cụm danh từ dùng để chỉ hình phạt tử hình.}

\vocabentry{Conviction}
{/kənˈvɪkʃn/ (n)}
{A formal declaration that someone is guilty of a criminal offense.}
{He has several prior convictions for driving under the influence of alcohol.}
{Đây là một danh từ dùng để chỉ sự kết án hoặc bản án có tội.}

\vocabentry{Custody}
{/ˈkʌstədi/ (n)}
{The protective care or guardianship of someone or something; parental responsibility.}
{The suspect is currently in police custody awaiting his first court appearance.}
{Trong luật pháp, đây là danh từ chỉ việc bị tạm giam hoặc quyền nuôi con.}

\vocabentry{Defendant}
{/dɪˈfendənt/ (n)}
{An individual, company, or institution sued or accused in a court of law.}
{The defendant maintained his innocence throughout the entire trial.}
{Đây là một danh từ dùng để chỉ bị cáo trong một vụ án.}

\vocabentry{Deterrence}
{/dɪˈterəns/ (n)}
{The action of discouraging an action or event through instilling doubt or fear of the consequences.}
{Long prison sentences are often used as a means of deterrence against serious crimes.}
{Đây là một danh từ dùng để chỉ sự răn đe, ngăn chặn tội phạm.}

\vocabentry{Embezzlement}
{/ɪmˈbezlmənt/ (n)}
{Theft or misappropriation of funds placed in one's trust or belonging to one's employer.}
{The accountant was charged with embezzlement after stealing thousands of dollars from the firm.}
{Đây là một danh từ dùng để chỉ tội tham ô tài sản.}

\vocabentry{Felony}
{/ˈfeləni/ (n)}
{A crime, typically one involving violence, regarded as more serious than a misdemeanor.}
{Armed robbery is classified as a felony and carries a much harsher sentence.}
{Đây là một danh từ dùng để chỉ trọng tội, tội đặc biệt nghiêm trọng.}

\vocabentry{Forgery}
{/ˈfɔːrdʒəri/ (n)}
{The action of forging or producing a copy of a document, signature, banknote, or work of art.}
{He was arrested for forgery after trying to use a fake passport at the airport.}
{Đây là một danh từ dùng để chỉ tội làm giả giấy tờ hoặc chữ ký.}

\vocabentry{Fraud}
{/frɔːd/ (n)}
{Wrongful or criminal deception intended to result in financial or personal gain.}
{Credit card fraud has become increasingly common in the era of online shopping.}
{Đây là một danh từ dùng để chỉ tội lừa đảo, thường liên quan đến tài chính.}

\vocabentry{Homicide}
{/ˈhɑːmɪsaɪd/ (n)}
{The killing of one person by another.}
{The detectives are treating the death as a homicide and are looking for suspects.}
{Đây là một danh từ trang trọng dùng để chỉ tội giết người.}

\vocabentry{Imprisonment}
{/ɪmˈprɪznmənt/ (n)}
{The state of being imprisoned; captivity.}
{The judge sentenced the criminal to ten years of imprisonment for his crimes.}
{Đây là một danh từ dùng để chỉ sự bỏ tù hoặc hình phạt tù túng.}

\vocabentry{Incarceration}
{/ɪnˌkɑːrsəˈreɪʃn/ (n)}
{The state of being confined in prison.}
{High rates of incarceration in the country have led to overcrowding in prisons.}
{Đây là một danh từ trang trọng đồng nghĩa với imprisonment, chỉ việc bị giam giữ.}

\vocabentry{Indictment}
{/ɪnˈdaɪtmənt/ (n)}
{A formal charge or accusation of a serious crime.}
{The grand jury issued an indictment against the governor for corruption.}
{Đây là một danh từ dùng để chỉ bản cáo trạng chính thức của tòa án.}

\vocabentry{Juvenile delinquency}
{/ˌdʒuːvənl dɪˈlɪŋkwənsi/ (n)}
{Habitual antisocial or unlawful behavior by young people.}
{Poverty and lack of education are often cited as major causes of juvenile delinquency.}
{Đây là một cụm danh từ dùng để chỉ tình trạng tội phạm vị thành niên.}

\vocabentry{Larceny}
{/ˈlɑːrsəni/ (n)}
{Theft of personal property.}
{He was charged with grand larceny after stealing an expensive sports car.}
{Đây là một danh từ pháp lý dùng để chỉ tội trộm cắp tài sản cá nhân.}

\vocabentry{Manslaughter}
{/ˈmænslɔːtər/ (n)}
{The crime of killing a human being without malice aforethought.}
{The driver was charged with involuntary manslaughter after the fatal accident.}
{Đây là một danh từ dùng để chỉ tội ngộ sát, giết người không cố ý.}

\vocabentry{Misdemeanor}
{/ˌmɪsdɪˈmiːnər/ (n)}
{A minor wrongdoing.}
{Shoplifting is often classified as a misdemeanor for first-time offenders.}
{Đây là một danh từ dùng để chỉ tội ít nghiêm trọng hơn trọng tội, thường là vi phạm hành chính.}

\vocabentry{Money laundering}
{/ˈmʌni ˌlɔːndərɪŋ/ (n)}
{The concealment of the origins of illegally obtained money.}
{Organized crime groups use shell companies for money laundering and tax evasion.}
{Đây là một cụm danh từ dùng để chỉ hành vi rửa tiền.}

\vocabentry{Pardon}
{/ˈpɑːrdn/ (n/v)}
{The action of forgiving or being forgiven for an error or offense.}
{The President granted a presidential pardon to the former soldier.}
{Trong luật pháp, đây là sự ân xá hoặc tha bổng.}

\vocabentry{Parole}
{/pəˈroʊl/ (n)}
{The release of a prisoner temporarily or permanently before the completion of a sentence.}
{He was released on parole after serving five years of his eight-year sentence.}
{Đây là một danh từ dùng để chỉ sự phóng thích có điều kiện/được tha tù trước thời hạn.}

\vocabentry{Perjury}
{/ˈpɜːrdʒəri/ (n)}
{The offense of willfully telling an untruth in a court after having taken an oath.}
{The witness was charged with perjury after it was proven that she lied under oath.}
{Đây là một danh từ dùng để chỉ tội khai man trước tòa.}

\vocabentry{Plaintiff}
{/ˈpleɪntɪf/ (n)}
{A person who brings a case against another in a court of law.}
{The plaintiff is seeking damages for the injuries he sustained in the accident.}
{Đây là một danh từ dùng để chỉ nguyên đơn, người đi kiện.}

\vocabentry{Prosecutor}
{/ˈprɑːsɪkjuːtər/ (n)}
{A person, especially a public official, who institutes legal proceedings against someone.}
{The public prosecutor presented compelling evidence that linked the suspect to the crime.}
{Đây là một danh từ dùng để chỉ công tố viên hoặc kiểm sát viên.}

\vocabentry{Recidivism}
{/rɪˈsɪdɪvɪzəm/ (n)}
{The tendency of a convicted criminal to reoffend.}
{Rehabilitation programs in prisons aim to reduce the rate of recidivism among ex-convicts.}
{Đây là một danh từ dùng để chỉ sự tái phạm tội.}

\vocabentry{Rehabilitation}
{/ˌriːəˌbɪlɪˈteɪʃn/ (n)}
{The action of restoring someone to health or normal life through training and therapy after imprisonment.}
{The prison focus should shift from punishment to the rehabilitation of offenders.}
{Đây là một danh từ dùng để chỉ sự cải tạo hoặc phục hồi nhân phẩm.}

\vocabentry{Sentence}
{/ˈsentəns/ (n/v)}
{The punishment assigned to a defendant found guilty by a court.}
{The judge handed down a life sentence for the horrific crime.}
{Đây là một danh từ/động từ dùng để chỉ bản án hoặc tuyên án.}

\vocabentry{Shoplifting}
{/ˈʃɑːplɪftɪŋ/ (n)}
{The criminal action of stealing goods from a shop while pretending to be a customer.}
{Shoplifting remains a major problem for retailers, costing them millions each year.}
{Đây là một danh từ dùng để chỉ tội trộm cắp ở cửa hàng.}

\vocabentry{Smuggling}
{/ˈsmʌɡlɪŋ/ (n)}
{The illegal movement of goods into or out of a country.}
{The border patrol arrested three men for smuggling illegal drugs across the border.}
{Đây là một danh từ dùng để chỉ hành vi buôn lậu.}

\vocabentry{Solicitor}
{/səˈlɪsɪtər/ (n)}
{A member of the legal profession qualified to deal with conveyancing, the drawing up of wills, etc.}
{You should consult a solicitor before signing any legal documents.}
{Đây là một danh từ dùng để chỉ luật sư tư vấn hoặc trạng sư.}

\vocabentry{Suspect}
{/ˈsʌspekt/ (n/v)}
{A person thought to be guilty of a crime or offense.}
{The police have identified a prime suspect in the ongoing murder investigation.}
{Đây là một danh từ/động từ dùng để chỉ nghi phạm hoặc nghi ngờ ai đó.}

\vocabentry{Testimony}
{/ˈtestɪmoʊni/ (n)}
{A formal written or spoken statement, especially one given in a court of law.}
{The jury was moved by the emotional testimony provided by the victim's mother.}
{Đây là một danh từ dùng để chỉ lời khai hoặc chứng cứ tại tòa.}

\vocabentry{Theft}
{/θeft/ (n)}
{The action or crime of stealing.}
{Identity theft is a growing concern for many people who use the internet for banking.}
{Đây là một danh từ dùng để chỉ tội trộm cắp nói chung.}

\vocabentry{Trial}
{/ˈtraɪəl/ (n)}
{A formal examination of evidence before a judge, and typically before a jury.}
{The high-profile murder trial is expected to last for several months.}
{Đây là một danh từ dùng để chỉ vụ xét xử hoặc phiên tòa.}

\vocabentry{Verdict}
{/ˈvɜːrdɪkt/ (n)}
{A decision on a disputed issue in a civil or criminal case.}
{The jury took only two hours to reach a unanimous verdict of guilty.}
{Đây là một danh từ dùng để chỉ phán quyết của bồi thẩm đoàn hoặc tòa án.}

\vocabentry{Witness}
{/ˈwɪtnəs/ (n/v)}
{A person who sees an event, typically a crime or accident, take place.}
{Several witnesses came forward to provide information about the hit-and-run accident.}
{Đây là một danh từ/động từ dùng để chỉ nhân chứng hoặc chứng kiến.}

\vocabentry{Bail}
{/beɪl/ (n/v)}
{The temporary release of an accused person awaiting trial, sometimes on condition that a sum of money is lodged to guarantee their appearance in court.}
{The judge set bail at \$50,000 for the defendant.}
{Đây là một danh từ/động từ dùng để chỉ tiền bảo lãnh hoặc việc bảo lãnh.}

\vocabentry{Bribery}
{/ˈbraɪbəri/ (n)}
{The giving or offering of a bribe.}
{Corruption and bribery are major obstacles to economic development in some countries.}
{Lặp lại để nhấn mạnh nghĩa danh từ chỉ hành vi hối lộ.}

\vocabentry{Criminology}
{/ˌkrɪmɪˈnɑːlədʒi/ (n)}
{The scientific study of crime and criminals.}
{She decided to major in criminology to understand the root causes of criminal behavior.}
{Đây là một danh từ dùng để chỉ tội phạm học.}

\vocabentry{Cybercrime}
{/ˈsaɪbərkraɪm/ (n)}
{Criminal activities carried out by means of computers or the internet.}
{The government is investing more resources to combat the rise in global cybercrime.}
{Đây là một danh từ dùng để chỉ tội phạm mạng.}

\vocabentry{Evidence}
{/ˈevɪdəns/ (n)}
{The available body of facts or information indicating whether a belief or proposition is true or valid.}
{DNA evidence played a crucial role in exonerating the man who was wrongly convicted.}
{Đây là một danh từ dùng để chỉ bằng chứng/chứng cứ.}

\vocabentry{Extortion}
{/ɪkˈstɔːrʃn/ (n)}
{The practice of obtaining something, especially money, through force or threats.}
{The gang members were arrested for extortion after threatening local business owners.}
{Đây là một danh từ dùng để chỉ tội tống tiền.}

\vocabentry{Heist}
{/haɪst/ (n)}
{A robbery. (Đây là một danh từ (không trang trọng) dùng để chỉ vụ cướp tài sản lớn, thường là từ ngân hàng hoặc bảo tàng.).}
{The movie is about a group of professional thieves planning a high-stakes diamond heist.}
{}

\vocabentry{Hijacking}
{/ˈhaɪdʒækɪŋ/ (n)}
{Illegally seizing (an aircraft, ship, or vehicle) in transit and forcing it to go to a different destination.}
{Security at airports was tightened significantly following the hijacking incident.}
{Đây là một danh từ dùng để chỉ hành vi không tặc hoặc cướp phương tiện.}

\vocabentry{Infraction}
{/ɪnˈfrækʃn/ (n)}
{A violation or infringement of a law or agreement.}
{A parking ticket is a minor infraction that results in a fine rather than a prison sentence.}
{Đây là một danh từ dùng để chỉ sự vi phạm, thường là những lỗi nhẹ.}

\vocabentry{Investigation}
{/ɪnˌvestɪˈɡeɪʃn/ (n)}
{The action of investigating something or someone; formal or systematic examination or research.}
{The FBI has launched a full investigation into the massive data breach.}
{Đây là một danh từ dùng để chỉ cuộc điều tra.}

\vocabentry{Justice}
{/ˈdʒʌstɪs/ (n)}
{Just behavior or treatment.}
{Many people took to the streets to demand justice for the victims of the tragedy.}
{Đây là một danh từ dùng để chỉ sự công lý hoặc hệ thống tư pháp.}

\vocabentry{Kidnapping}
{/ˈkɪdnæpɪŋ/ (n)}
{The action of abducting someone and holding them captive, typically to obtain a ransom.}
{The kidnapping of the businessman sent shockwaves throughout the local community.}
{Đây là một danh từ dùng để chỉ vụ bắt cóc tống tiền.}

\vocabentry{Lawyer}
{/ˈlɔɪər/ (n)}
{A person who practices or studies law; an attorney or a counselor.}
{If you are arrested, you have the right to speak to a lawyer before answering any questions.}
{Đây là một danh từ dùng để chỉ luật sư nói chung.}

\vocabentry{Legal}
{/ˈliːɡl/ (adj)}
{Relating to the law.}
{You should seek professional legal advice before entering into any business contract.}
{Đây là một tính từ dùng để chỉ những gì thuộc về luật pháp hoặc hợp pháp.}

\vocabentry{Looting}
{/ˈluːtɪŋ/ (n)}
{Steal goods from (a place), typically during a war or riot.}
{Widespread looting was reported in the city following the devastating earthquake.}
{Đây là một danh từ dùng để chỉ hành vi hôi của hoặc cướp bóc trong lúc loạn lạc.}

\vocabentry{Malice}
{/ˈmælɪs/ (n)}
{The intention or desire to do evil; ill will.}
{To prove murder, the prosecutor must show that the defendant acted with malice.}
{Đây là một danh từ dùng để chỉ ác ý hoặc ý đồ xấu, quan trọng trong việc định tội.}

\vocabentry{Murder}
{/ˈmɜːrdər/ (n/v)}
{The unlawful premeditated killing of one human being by another.}
{He was sentenced to life in prison for the murder of his business partner.}
{Đây là một danh từ/động từ dùng để chỉ tội giết người có mưu sát.}

\vocabentry{Offender}
{/əˈfendər/ (n)}
{A person who commits an illegal act.}
{The program aims to reintegrate young offenders back into society through vocational training.}
{Đây là một danh từ dùng để chỉ người phạm tội/kẻ phạm tội.}

\vocabentry{Penalty}
{/ˈpenəlti/ (n)}
{A punishment imposed for breaking a law, rule, or contract.}
{There are severe penalties for anyone caught driving under the influence of drugs.}
{Đây là một danh từ dùng để chỉ hình phạt hoặc tiền phạt.}

\vocabentry{Plaintiff}
{/ˈpleɪntɪf/ (n)}
{A person who brings a case against another in a court of law.}
{The plaintiff's lawyer argued that the company was responsible for the defective product.}
{Lặp lại để ghi nhớ nghĩa nguyên đơn.}

\vocabentry{Police}
{/pəˈliːs/ (n)}
{The civil force of a national or local government, responsible for the prevention and detection of crime and the maintenance of public order.}
{The police were called to the scene of the domestic disturbance.}
{Đây là một danh từ dùng để chỉ lực lượng cảnh sát.}

\vocabentry{Prison}
{/ˈprɪzn/ (n)}
{A building in which people are legally held as a punishment for a crime they have committed or while awaiting trial.}
{The overcrowding in state prisons is a major concern for human rights organizations.}
{Đây là một danh từ dùng để chỉ nhà tù.}

\vocabentry{Probation}
{/proʊˈbeɪʃn/ (n)}
{The release of an offender from detention, subject to a period of good behavior under supervision.}
{Instead of being sent to prison, the teenager was placed on two years of probation.}
{Đây là một danh từ dùng để chỉ án treo hoặc thời gian thử thách.}

\vocabentry{Punishment}
{/ˈpʌnɪʃmənt/ (n)}
{The infliction or imposition of a penalty as retribution for an offense.}
{Many believe that the punishment should always fit the crime.}
{Đây là một danh từ dùng để chỉ sự trừng phạt.}

\vocabentry{Ransom}
{/ˈrænsəm/ (n/v)}
{A sum of money demanded or paid for the release of a captive.}
{The kidnappers demanded a five-million dollar ransom for the safe return of the girl.}
{Đây là một danh từ/động từ dùng để chỉ tiền chuộc hoặc đòi tiền chuộc.}

\vocabentry{Rehabilitation}
{/ˌriːəˌbɪlɪˈteɪʃn/ (n)}
{The action of restoring someone to health or normal life through training and therapy.}
{Effective rehabilitation can help reduce the chances of a criminal reoffending.}
{Lặp lại để nhấn mạnh nghĩa sự cải tạo.}

\vocabentry{Release}
{/rɪˈliːs/ (n/v)}
{Allow or enable to escape from confinement; set free.}
{The innocent man was finally released from prison after twenty years.}
{Đây là một danh từ/động từ dùng để chỉ sự phóng thích hoặc thả tự do.}

\vocabentry{Retribution}
{/ˌretrɪˈbjuːʃn/ (n)}
{Punishment inflicted on someone as vengeance for a wrong or criminal act.}
{Some victims' families seek retribution, while others focus on finding closure.}
{Đây là một danh từ dùng để chỉ sự báo thù hoặc sự trừng phạt thích đáng.}

\vocabentry{Robbery}
{/ˈrɑːbəri/ (n)}
{The action of taking property unlawfully from a person or place by force or threat of force.}
{An armed robbery took place at the jewelry store during broad daylight.}
{Đây là một danh từ dùng để chỉ tội cướp tài sản dùng vũ lực.}

\vocabentry{Smuggling}
{/ˈsmʌɡlɪŋ/ (n)}
{The illegal movement of goods into or out of a country.}
{The police are working hard to crack down on the smuggling of counterfeit goods.}
{Lặp lại để nhấn mạnh nghĩa buôn lậu.}

\vocabentry{Solicitor}
{/səˈlɪsɪtər/ (n)}
{A member of the legal profession.}
{You should contact a solicitor to help you with the legal aspects of buying a house.}
{Lặp lại để nhấn mạnh nghĩa luật sư tư vấn.}

\vocabentry{Stalking}
{/ˈstɔːkɪŋ/ (n)}
{The act of following or harassing someone in an obsessive or threatening way.}
{She contacted the police after becoming a victim of online stalking.}
{Đây là một danh từ dùng để chỉ hành vi lén lút theo dõi/quấy rối người khác.}

\vocabentry{Statute}
{/ˈstætʃuːt/ (n)}
{A written law passed by a legislative body.}
{The protection of endangered species is enshrined in national statute.}
{Đây là một danh từ dùng để chỉ đạo luật hoặc quy chế.}

\vocabentry{Subpoena}
{/səˈpiːnə/ (n/v)}
{A writ ordering a person to attend a court.}
{The defense lawyer requested a subpoena for several key witnesses.}
{Đây là một danh từ/động từ dùng để chỉ lệnh hầu tòa hoặc trát tòa.}

\vocabentry{Surveillance}
{/sɜːrˈveɪləns/ (n)}
{Close observation, especially of a suspected spy or criminal.}
{The suspect was under 24-hour police surveillance for several weeks.}
{Đây là một danh từ dùng để chỉ sự giám sát/theo dõi.}

\vocabentry{Suspect}
{/ˈsʌspekt/ (n/v)}
{A person thought to be guilty of a crime.}
{The police have arrested a suspect in connection with the bank robbery.}
{Lặp lại để ghi nhớ nghĩa nghi phạm.}

\vocabentry{Theft}
{/θeft/ (n)}
{The action or crime of stealing.}
{He was caught on camera committing a petty theft at the local convenience store.}
{Lặp lại để ghi nhớ nghĩa tội trộm cắp.}

\vocabentry{Trial}
{/ˈtraɪəl/ (n)}
{A formal examination of evidence before a judge.}
{The trial had to be postponed because a key witness was unable to attend.}
{Lặp lại để ghi nhớ nghĩa vụ xét xử.}

\vocabentry{Vandalism}
{/ˈvændəlɪzəm/ (n)}
{Action involving deliberate destruction of or damage to public or private property.}
{The local park has suffered from increasing acts of vandalism recently.}
{Đây là một danh từ dùng để chỉ hành vi phá hoại tài sản công cộng.}

\vocabentry{Verdict}
{/ˈvɜːrdɪkt/ (n)}
{A decision on a disputed issue in a civil or criminal case.}
{The judge will deliver his verdict on the case next Friday morning.}
{Lặp lại để ghi nhớ nghĩa phán quyết.}

\vocabentry{Victim}
{/ˈvɪktɪm/ (n)}
{A person harmed, injured, or killed as a result of a crime, accident, or other event or action.}
{Support groups provide essential services for the victims of domestic violence.}
{Đây là một danh từ dùng để chỉ nạn nhân.}

\vocabentry{Violation}
{/ˌvaɪəˈleɪʃn/ (n)}
{The action of violating someone or something.}
{The company was fined for multiple violations of health and safety regulations.}
{Đây là một danh từ dùng để chỉ sự vi phạm luật pháp hoặc quyền lợi.}

\vocabentry{Warrant}
{/ˈwɔːrənt/ (n)}
{A document issued by a legal or government official authorizing the police or some other body to make an arrest, search premises.}
{The police obtained a search warrant before entering the suspect's house.}
{Đây là một danh từ dùng để chỉ lệnh bắt giữ hoặc lệnh khám xét.}

\vocabentry{White-collar crime}
{/ˌwaɪt ˈkɑːlər kraɪm/ (n)}
{Financially motivated nonviolent crime committed by business and government professionals.}
{Insider trading and tax evasion are typical examples of white-collar crime.}
{Đây là một cụm danh từ dùng để chỉ tội phạm cổ cồn trắng/tội phạm kinh tế trí thức.}

\vocabentry{Witness}
{/ˈwɪtnəs/ (n/v)}
{A person who sees an event take place.}
{The witness provided a detailed description of the getaway car to the police.}
{Lặp lại để ghi nhớ nghĩa nhân chứng.}

\vocabentry{Accusation}
{/ˌækjuˈzeɪʃn/ (n)}
{A charge or claim that someone has done something illegal or wrong.}
{He denied the accusations of corruption made against him by the media.}
{Đây là một danh từ dùng để chỉ sự cáo buộc.}

\vocabentry{Alibi}
{/ˈæləbaɪ/ (n)}
{A claim or piece of evidence that one was elsewhere when an act, typically a criminal one, is alleged to have taken place.}
{The suspect had a solid alibi; he was at work when the crime occurred.}
{Đây là một danh từ dùng để chỉ chứng cứ ngoại phạm.}

\vocabentry{Allegation}
{/ˌæləˈɡeɪʃn/ (n)}
{A claim or assertion that someone has done something illegal or wrong, typically one made without proof.}
{The police are investigating allegations of physical abuse at the nursing home.}
{Đây là một danh từ dùng để chỉ sự khẳng định/cáo buộc chưa có bằng chứng.}

\vocabentry{Appeal}
{/əˈpiːl/ (n/v)}
{Apply to a higher court for a reversal of the decision of a lower court.}
{The defendant decided to appeal against his conviction for armed robbery.}
{Đây là một danh từ/động từ dùng để chỉ sự kháng cáo.}

\vocabentry{Appellate}
{/əˈpelət/ (adj)}
{(Of a court) concerned with or dealing with applications for decisions to be reversed.}
{The appellate court overturned the lower court's ruling due to a legal error.}
{Đây là một tính từ dùng để chỉ tòa án phúc thẩm.}

\vocabentry{Armed robbery}
{/ˌɑːrmd ˈrɑːbəri/ (n)}
{The crime of taking property from a person or place using a weapon.}
{The shopkeeper was threatened during the armed robbery but was not physically harmed.}
{Đây là một cụm danh từ dùng để chỉ vụ cướp có vũ trang.}

\vocabentry{Attorney}
{/əˈtɜːrni/ (n)}
{A person, typically a lawyer, appointed to act for another in business or legal matters.}
{You have the right to an attorney; if you cannot afford one, the state will provide one.}
{Đây là một danh từ chỉ luật sư ủy quyền/biện hộ.}

\vocabentry{Ban}
{/bæn/ (n/v)}
{Officially or legally prohibit.}
{The government has decided to ban the sale of tobacco to anyone under 21.}
{Đây là một danh từ/động từ dùng để chỉ lệnh cấm hoặc việc cấm đoán.}

\vocabentry{Barrister}
{/ˈbærɪstər/ (n)}
{A lawyer entitled to practice as an advocate, particularly in the higher courts.}
{The barrister delivered a powerful closing argument that swayed the jury.}
{Đây là một danh từ dùng để chỉ luật sư tranh tụng trước tòa.}

\vocabentry{Blackmail}
{/ˈblækmeɪl/ (n/v)}
{The action, treated as a criminal offense, of demanding payment or another benefit from someone in return for not revealing compromising information about them.}
{He was arrested for attempting to blackmail a famous Hollywood actor.}
{Đây là một danh từ/động từ dùng để chỉ hành vi tống tiền bằng bí mật cá nhân.}

\vocabentry{Bootlegging}
{/ˈbuːtleɡɪŋ/ (n)}
{Make, distribute, or sell (illicit goods, especially liquor, computer software, or recordings) illegally.}
{Bootlegging of alcohol was a major problem during the Prohibition era.}
{Đây là một danh từ dùng để chỉ hành vi sản xuất/buôn bán hàng cấm, lậu.}

\vocabentry{Bounty}
{/ˈbaʊnti/ (n)}
{A sum paid for killing or capturing a person or animal.}
{The government has offered a large bounty for any information leading to the arrest of the terrorist.}
{Đây là một danh từ dùng để chỉ tiền thưởng cho việc bắt tội phạm.}

\vocabentry{Breathalyzer}
{/ˈbreθəlaɪzər/ (n)}
{A device used by police for measuring the amount of alcohol in a driver's breath.}
{The police officer asked the driver to blow into the breathalyzer.}
{Đây là một danh từ dùng để chỉ máy đo nồng độ cồn.}

\vocabentry{Bribe}
{/braɪb/ (n/v)}
{Persuade (someone) to act in one's favor, typically illegally or dishonestly, by a gift of money or other inducement.}
{He tried to bribe the border official to let him through without a visa.}
{Đây là một danh từ/động từ dùng để chỉ của hối lộ hoặc việc đút lót.}

\vocabentry{Case}
{/keɪs/ (n)}
{An instance of a particular situation; an incident or series of incidents being investigated by the police or heard in a court.}
{The detective spent months working on the cold case that had remained unsolved for years.}
{Đây là một danh từ dùng để chỉ vụ án hoặc hồ sơ pháp lý.}

\vocabentry{Cell}
{/sel/ (n)}
{A small room in which a prisoner is locked.}
{The prisoner was confined to his cell for twenty-three hours a day.}
{Đây là một danh từ dùng để chỉ buồng giam.}

\vocabentry{Charge}
{/tʃɑːrdʒ/ (n/v)}
{Accuse someone of something, especially an offense under law.}
{The police have decided to drop all charges against the suspect.}
{Đây là một danh từ/động từ dùng để chỉ sự cáo buộc hoặc buộc tội.}

\vocabentry{Civil law}
{/ˌsɪvl ˈlɔː/ (n)}
{The system of law concerned with private relations between members of a community rather than criminal, military, or religious affairs.}
{In civil law cases, the burden of proof is lower than in criminal cases.}
{Đây là một cụm danh từ dùng để chỉ luật dân sự.}

\vocabentry{Commute}
{/kəˈmjuːt/ (v)}
{Reduce (a judicial sentence, especially a sentence of death) to one less severe.}
{The governor decided to commute the prisoner's death sentence to life in prison.}
{Trong tư pháp, đây là việc giảm án.}

\vocabentry{Compensation}
{/ˌkɑːmpenˈseɪʃn/ (n)}
{Something, typically money, awarded to someone as a recompense for loss, injury, or suffering.}
{The victim is seeking compensation for the medical expenses and lost wages.}
{Đây là một danh từ dùng để chỉ sự bồi thường.}

\vocabentry{Confession}
{/kənˈfeʃn/ (n)}
{A formal statement admitting that one is guilty of a crime.}
{The police obtained a full confession from the suspect during the interrogation.}
{Đây là một danh từ dùng để chỉ lời thú tội.}

\vocabentry{Confiscate}
{/ˈkɑːnfɪskeɪt/ (v)}
{Take or seize (someone's property) with authority.}
{Customs officials confiscated the illegal ivory that was hidden in the luggage.}
{Đây là một động từ dùng để chỉ việc tịch thu tài sản.}

\vocabentry{Consensus}
{/kənˈsensəs/ (n)}
{A general agreement.}
{The jury was unable to reach a consensus, resulting in a hung jury.}
{Trong bồi thẩm đoàn, đây là sự đồng thuận về phán quyết.}

\vocabentry{Contraband}
{/ˈkɑːntrəbænd/ (n)}
{Goods that have been imported or exported illegally.}
{The ship was searched for contraband drugs and weapons.}
{Đây là một danh từ dùng để chỉ hàng lậu/hàng cấm.}

\vocabentry{Coroner}
{/ˈkɔːrənər/ (n)}
{An official who investigates violent, sudden, or suspicious deaths.}
{The coroner's report confirmed that the cause of death was accidental drowning.}
{Đây là một danh từ dùng để chỉ nhân viên điều tra nguyên nhân cái chết/pháp y.}

\vocabentry{Counsel}
{/ˈkaʊnsl/ (n)}
{The lawyer or lawyers conducting a case.}
{The judge asked the defense counsel to proceed with their opening statement.}
{Đây là một danh từ đồng nghĩa với lawyer, thường dùng để chỉ luật sư bào chữa trong phiên tòa.}

\vocabentry{Counterfeit}
{/ˈkaʊntərfɪt/ (adj/v)}
{Made in exact imitation of something valuable with the intention to deceive or defraud.}
{The police seized thousands of dollars in counterfeit banknotes.}
{Đây là một tính từ/động từ dùng để chỉ hàng giả hoặc việc làm giả.}

\vocabentry{Court}
{/kɔːrt/ (n)}
{A tribunal presided over by a judge, judges, or a magistrate in civil and criminal cases.}
{The case will be heard in the high court next Monday.}
{Đây là một danh từ dùng để chỉ tòa án.}

\vocabentry{Crime}
{/kraɪm/ (n)}
{An action or omission that constitutes an offense that may be prosecuted by the state and is punishable by law.}
{Violent crime rates have been decreasing in the city over the last decade.}
{Đây là một danh từ dùng để chỉ tội phạm hoặc hành vi phạm tội.}

\vocabentry{Criminal}
{/ˈkrɪmɪnl/ (n/adj)}
{A person who has committed a crime.}
{The dangerous criminal was captured after a high-speed car chase.}
{Đây là một danh từ/tính từ dùng để chỉ tội phạm hoặc có tính chất phạm tội.}

\vocabentry{Cross-examination}
{/ˌkrɔːs ɪɡˌzæmɪˈneɪʃn/ (n)}
{The formal interrogation of a witness called by the other party in a court of law to challenge or extend testimony already given.}
{During cross-examination, the defense lawyer pointed out several inconsistencies in the witness's story.}
{Đây là một danh từ dùng để chỉ sự thẩm vấn chéo tại tòa.}

\vocabentry{Damages}
{/ˈdæmɪdʒɪz/ (n)}
{A sum of money claimed or awarded in compensation for a loss or an injury.}
{The court awarded the plaintiff \$100,000 in damages for pain and suffering.}
{Trong luật dân sự, đây là khoản tiền bồi thường thiệt hại.}

\vocabentry{Decree}
{/dɪˈkriː/ (n/v)}
{An official order issued by a legal authority.}
{The judge issued a decree for the immediate release of the prisoner.}
{Đây là một danh từ/động từ dùng để chỉ sắc lệnh hoặc việc ra lệnh.}

\vocabentry{Deliberation}
{/dɪˌlɪbəˈreɪʃn/ (n)}
{Long and careful consideration or discussion.}
{After three days of deliberation, the jury finally reached a decision.}
{Trong tòa án, đây là sự thảo luận kín của bồi thẩm đoàn để đưa ra phán quyết.}

\vocabentry{Deposition}
{/ˌdepəˈzɪʃn/ (n)}
{The process of giving sworn evidence.}
{The lawyer took a deposition from the doctor regarding the victim's injuries.}
{Đây là một danh từ dùng để chỉ việc lấy lời khai có tuyên thệ bên ngoài tòa án.}

\vocabentry{Detective}
{/dɪˈtektɪv/ (n)}
{A person, especially a police officer, whose occupation is to investigate and solve crimes.}
{The lead detective spent years tracking down the serial killer.}
{Đây là một danh từ dùng để chỉ thám tử hoặc điều tra viên.}

\vocabentry{Detention}
{/dɪˈtenʃn/ (n)}
{The action of detaining someone or the state of being detained in official custody.}
{The suspect was held in detention for 48 hours without being charged.}
{Đây là một danh từ dùng để chỉ sự giam cầm hoặc tạm giữ.}

\vocabentry{Disorderly conduct}
{/dɪsˌɔːrdərli ˈkɑːndʌkt/ (n)}
{A legal term used to describe behavior that is offensive, annoying, or disruptive to others.}
{He was arrested for disorderly conduct after getting into a fight at the bar.}
{Đây là một cụm danh từ dùng để chỉ hành vi gây rối trật tự công cộng.}

\vocabentry{Docket}
{/ˈdɑːkɪt/ (n)}
{A calendar or list of cases for trial or people having cases pending.}
{The court's docket is full for the next six months, leading to long delays.}
{Đây là một danh từ dùng để chỉ danh sách các vụ án sắp xét xử.}

\vocabentry{Domestic violence}
{/dəˌmestɪk ˈvaɪələns/ (n)}
{Violent or aggressive behavior within the home, typically involving the violent abuse of a spouse or partner.}
{The charity provides a safe haven for women and children who are fleeing domestic violence.}
{Đây là một cụm danh từ dùng để chỉ bạo lực gia đình.}

\vocabentry{Drug trafficking}
{/ˈdrʌɡ ˌtræfɪkɪŋ/ (n)}
{The production, distribution, and sale of illegal drugs.}
{Drug trafficking is a major source of income for many international criminal organizations.}
{Đây là một cụm danh từ dùng để chỉ tội buôn lậu ma túy.}

\vocabentry{Drunken driving}
{/ˌdrʌŋkən ˈdraɪvɪŋ/ (n)}
{The crime of driving a vehicle with excess alcohol in the blood.}
{Drunken driving is a serious offense that can lead to fatal accidents.}
{Đây là một cụm danh từ dùng để chỉ hành vi lái xe khi say rượu.}

\vocabentry{Due process}
{/ˌduː ˈproʊses/ (n)}
{Fair treatment through the normal judicial system, especially as a citizen's entitlement.}
{The legal system must ensure that every citizen receives due process of law.}
{Đây là một cụm danh từ dùng để chỉ quy trình tố tụng công bằng.}

\vocabentry{Embezzle}
{/ɪmˈbezl/ (v)}
{Steal or misappropriate (money placed in one's trust or belonging to one's employer).}
{He was accused of trying to embezzle funds from the pension scheme.}
{Động từ của embezzlement, chỉ việc thụ thụt/tham ô.}

\vocabentry{Exonerate}
{/ɪɡˈzɑːnəreɪt/ (v)}
{Absolve (someone) from blame for a fault or wrongdoing, especially after due consideration of the case.}
{New DNA testing has helped to exonerate many people who were wrongly convicted.}
{Đây là một động từ dùng để chỉ việc giải oan hoặc minh oan.}

\vocabentry{Extradition}
{/ˌekstrəˈdɪʃn/ (n)}
{The action of extraditing a person accused or convicted of a crime.}
{The government is seeking the extradition of the suspect from the neighboring country.}
{Đây là một danh từ dùng để chỉ việc dẫn độ tội phạm.}

\vocabentry{Eye-witness}
{/ˈaɪ ˌwɪtnəs/ (n)}
{A person who has personally seen something happen and so can give a first-hand description of it.}
{The prosecution called an eye-witness who saw the defendant leaving the crime scene.}
{Đây là một danh từ dùng để chỉ nhân chứng trực tiếp kiến chứng sự việc.}

\vocabentry{Felon}
{/ˈfelən/ (n)}
{A person who has been convicted of a felony.}
{As a convicted felon, he is no longer allowed to own a firearm.}
{Đây là một danh từ dùng để chỉ tội phạm phạm trọng tội.}

\vocabentry{Fine}
{/faɪn/ (n/v)}
{A sum of money exacted as a penalty by a court of law or other authority.}
{He had to pay a heavy fine for speeding in a residential area.}
{Đây là một danh từ/động từ dùng để chỉ tiền phạt hoặc việc xử phạt bằng tiền.}

\vocabentry{Forensic}
{/fəˈrenzɪk/ (adj)}
{Relating to or denoting the application of scientific methods and techniques to the investigation of crime.}
{Forensic evidence like fingerprints and DNA helped to identify the killer.}
{Đây là một tính từ dùng để chỉ thuộc về pháp y/kỹ thuật hình sự.}

\vocabentry{Gallows}
{/ˈɡæloʊz/ (n)}
{A structure, typically of two uprights and a crosspiece, for the hanging of criminals.}
{In the past, many criminals met their end on the gallows in the town square.}
{Đây là một danh từ dùng để chỉ giá treo cổ.}

\vocabentry{Gang}
{/ɡæŋ/ (n)}
{An organized group of criminals.}
{The city has seen a rise in gang-related violence over the last year.}
{Đây là một danh từ dùng để chỉ băng đảng tội phạm.}

\vocabentry{Grand jury}
{/ˌɡrænd ˈdʒʊri/ (n)}
{A jury, normally of twenty-three jurors, selected to examine the validity of an accusation before trial.}
{The grand jury decided there was enough evidence to proceed with the trial.}
{Đây là một cụm danh từ dùng để chỉ đại bồi thẩm đoàn.}

\vocabentry{Guilty}
{/ˈɡɪlti/ (adj)}
{Culpable of or responsible for a specified wrongdoing.}
{The defendant pleaded guilty to the charges of theft and assault.}
{Đây là một tính từ dùng để chỉ việc có tội.}

\vocabentry{Hate crime}
{/ˈheɪt kraɪm/ (n)}
{A crime, typically one involving violence, that is motivated by prejudice on the basis of race, religion, sexual orientation, or other grounds.}
{The assault is being investigated as a hate crime because of the racial slurs used.}
{Đây là một cụm danh từ dùng để chỉ tội phạm do thù ghét/kỳ thị.}

\vocabentry{Hearing}
{/ˈhɪrɪŋ/ (n)}
{An opportunity to state one's case.}
{The preliminary hearing will determine if there is enough evidence for a trial.}
{Đây là một danh từ dùng để chỉ phiên điều trần hoặc phiên tòa sơ thẩm.}

\vocabentry{Homicide}
{/ˈhɑːmɪsaɪd/ (n)}
{The killing of one person by another.}
{The homicide rate in the country has reached its highest level in twenty years.}
{Lặp lại để nhấn mạnh nghĩa danh từ killing.}

\vocabentry{Hostage}
{/ˈhɑːstɪdʒ/ (n)}
{A person seized or held as security for the fulfillment of a condition.}
{The bank robbers took several employees as hostages during the standoff with police.}
{Đây là một danh từ dùng để chỉ con tin.}

\vocabentry{Hung jury}
{/ˌhʌŋ ˈdʒʊri/ (n)}
{A jury that is unable to reach a verdict of guilty or not guilty.}
{The judge declared a mistrial after the case resulted in a hung jury.}
{Đây là một cụm danh từ dùng để chỉ bồi thẩm đoàn không thống nhất được phán quyết.}

\vocabentry{Illegal}
{/ɪˈliːɡl/ (adj)}
{Contrary to or forbidden by law, especially criminal law.}
{It is illegal to drive a car without a valid driver's license.}
{Đây là một tính từ dùng để chỉ những gì bất hợp pháp.}

\vocabentry{Immunity}
{/ɪˈmjuːnəti/ (n)}
{Exemption from a service, duty, or liability.}
{The witness was granted immunity from prosecution in exchange for his testimony.}
{Trong luật pháp, đây là quyền miễn trừ.}

\vocabentry{Imprison}
{/ɪmˈprɪzn/ (v)}
{Put or keep in prison.}
{He was imprisoned for five years for his involvement in the political protests.}
{Động từ của prison, chỉ việc bỏ tù.}

\vocabentry{Incarcerate}
{/ɪnˈkɑːrsəreɪt/ (v)}
{Imprison or confine.}
{The government plans to incarcerate more drug dealers in high-security prisons.}
{Động từ trang trọng của imprison.}

\vocabentry{Innocent}
{/ˈɪnəsnt/ (adj)}
{Not guilty of a crime or offense.}
{The man spent years in prison before it was proven that he was completely innocent.}
{Đây là một tính từ dùng để chỉ sự vô tội.}

\vocabentry{Inmate}
{/ˈɪnmeɪt/ (n)}
{A person confined to an institution such as a prison or hospital.}
{The prison inmates are allowed to have visitors once a week.}
{Đây là một danh từ dùng để chỉ bạn tù hoặc người bị giam giữ.}

\vocabentry{Interrogation}
{/ɪnˌterəˈɡeɪʃn/ (n)}
{The action of interrogating someone.}
{The suspect was subjected to several hours of intense interrogation by the detectives.}
{Đây là một danh từ dùng để chỉ sự thẩm vấn hoặc tra hỏi.}

\vocabentry{Judge}
{/dʒʌdʒ/ (n/v)}
{A public official appointed to decide cases in a court of law.}
{The judge listened carefully to both sides of the argument before making a decision.}
{Đây là một danh từ/động từ dùng để chỉ thẩm phán hoặc việc xét xử.}

\vocabentry{Jury}
{/ˈdʒʊri/ (n)}
{A body of people (typically twelve in number) sworn to give a verdict in a legal case.}
{Members of the jury must be impartial and base their decision only on the evidence.}
{Đây là một danh từ dùng để chỉ bồi thẩm đoàn.}

\vocabentry{Justice}
{/ˈdʒʌstɪs/ (n)}
{Just behavior or treatment.}
{The quest for justice for the victims continues even after the trial has ended.}
{Lặp lại lần 2 để nhấn mạnh nghĩa sự công bằng/công lý.}

\vocabentry{Law}
{/lɔː/ (n)}
{The system of rules which a particular country or community recognizes as regulating the actions of its members.}
{No one is above the law, regardless of their status or wealth.}
{Đây là một danh từ dùng để chỉ luật pháp.}

\vocabentry{Lawsuit}
{/ˈlɔːsuːt/ (n)}
{A claim or dispute brought to a law court for adjudication.}
{The environmental group has filed a lawsuit against the factory for polluting the river.}
{Đây là một danh từ dùng để chỉ vụ kiện dân sự.}

\vocabentry{Legal aid}
{/ˌliːɡl ˈeɪd/ (n)}
{Payment from public funds allowed to people who cannot afford to pay for a lawyer.}
{People on low incomes can apply for legal aid to help with their court costs.}
{Đây là một cụm danh từ dùng để chỉ sự trợ giúp pháp lý miễn phí.}

\vocabentry{Life sentence}
{/ˌlaɪf ˈsentəns/ (n)}
{A sentence of imprisonment for the rest of a person's life.}
{He was given a life sentence for the double murder he committed.}
{Đây là một cụm danh từ dùng để chỉ bản án chung thân.}

\vocabentry{Litigation}
{/ˌlɪtɪˈɡeɪʃn/ (n)}
{The process of taking legal action.}
{The company is currently involved in a complex piece of litigation over patent rights.}
{Đây là một danh từ dùng để chỉ quá trình tranh tụng tại tòa.}

\vocabentry{Magistrate}
{/ˈmædʒɪstreɪt/ (n)}
{A civil officer or lay judge who administers the law.}
{The case was heard by a local magistrate who decided to impose a fine.}
{Đây là một danh từ dùng để chỉ thẩm phán sơ thẩm hoặc viên chức tư pháp.}

\vocabentry{Mandatory}
{/ˈmændətɔːri/ (adj)}
{Required by law or rules; compulsory.}
{There is a mandatory minimum sentence for certain drug-related offenses.}
{Đây là một tính từ dùng để chỉ những gì bắt buộc, ví dụ: "mandatory sentencing" - tuyên án bắt buộc.}

\vocabentry{Manslaughter}
{/ˈmænslɔːtər/ (n)}
{The crime of killing a human being without malice aforethought.}
{The charges were reduced from murder to manslaughter due to the circumstances.}
{Lặp lại để nhấn mạnh nghĩa tội ngộ sát.}

\vocabentry{Money laundering}
{/ˈmʌni ˌlɔːndərɪŋ/ (n)}
{The concealment of the origins of illegally obtained money.}
{The bank was fined for its failure to prevent money laundering through its accounts.}
{Lặp lại để ghi nhớ nghĩa tội rửa tiền.}

\vocabentry{Mugging}
{/ˈmʌɡɪŋ/ (n)}
{An act of attacking and robbing someone in a public place.}
{The tourist was traumatized after a mugging in a dark alleyway.}
{Đây là một danh từ dùng để chỉ vụ trấn lột/cướp giật trên đường phố.}

\vocabentry{Negligence}
{/ˈneɡlɪdʒəns/ (n)}
{Failure to take proper care in doing something.}
{The hospital was sued for medical negligence after the patient's condition worsened.}
{Đây là một danh từ dùng để chỉ sự sơ suất hoặc cẩu thả gây hậu quả pháp lý.}

\vocabentry{Not guilty}
{/ˌnɑːt ˈɡɪlti/ (adj)}
{Innocent of a crime.}
{The defendant pleaded not guilty and asked for a trial by jury.}
{Đây là một cụm tính từ dùng để chỉ sự vô tội.}

\vocabentry{Oath}
{/oʊθ/ (n)}
{A solemn promise, often invoking a divine witness, regarding one's future action or behavior.}
{All witnesses must take an oath to tell the truth before they give evidence.}
{Đây là một danh từ dùng để chỉ lời thề hoặc lời tuyên thệ tại tòa.}

\vocabentry{Offense}
{/əˈfens/ (n)}
{A breach of a law or rule; an illegal act.}
{This is his first offense, so the judge might be lenient with the sentence.}
{Đây là một danh từ dùng để chỉ hành vi phạm tội/lỗi vi phạm.}

\vocabentry{Order}
{/ˈɔːrdər/ (n/v)}
{An authoritative command, direction, or instruction.}
{The judge issued an order for the suspect to stay away from the victim.}
{Đây là một danh từ/động từ dùng để chỉ trật tự hoặc lệnh của tòa án.}

\vocabentry{Overrule}
{/ˌoʊvərˈruːl/ (v)}
{Reject or disallow by exercising one's superior authority. (Đây là một động từ dùng để chỉ việc bác bỏ (lời phản đối của luật sư) bởi thẩm phán.).}
{The judge overruled the prosecutor's objection and allowed the witness to continue.}
{}

\vocabentry{Pardon}
{/ˈpɑːrdn/ (n/v)}
{The action of forgiving an offense.}
{He was granted a full pardon after spent ten years in prison for a crime he didn't commit.}
{Lặp lại để ghi nhớ nghĩa ân xá.}

\vocabentry{Parole}
{/pəˈroʊl/ (n)}
{Release of a prisoner before completion of a sentence.}
{To stay on parole, he must meet with his parole officer every week.}
{Lặp lại lần 2 để nhấn mạnh nghĩa phóng thích có điều kiện.}

\vocabentry{Penal code}
{/ˌpiːnl ˈkoʊd/ (n)}
{A code of laws concerning crimes and offenses and their punishment.}
{The country is currently revising its penal code to include new types of cybercrime.}
{Đây là một cụm danh từ dùng để chỉ bộ luật hình sự.}

\vocabentry{Penalty}
{/ˈpenəlti/ (n)}
{A punishment imposed for breaking a law.}
{The maximum penalty for this crime is life imprisonment.}
{Lặp lại để ghi nhớ nghĩa hình phạt.}

\vocabentry{Perjury}
{/ˈpɜːrdʒəri/ (n)}
{The offense of willfully telling an untruth in a court.}
{He committed perjury when he lied about his whereabouts on the night of the murder.}
{Lặp lại lần 2 để nhấn mạnh tội khai man.}

\vocabentry{Plea}
{/pliː/ (n)}
{A formal statement by or on behalf of a defendant or prisoner, stating guilt or innocence.}
{The defendant entered a plea of not guilty to all of the charges.}
{Đây là một danh từ dùng để chỉ lời bào chữa hoặc sự tự nhận tội/vô tội.}

\vocabentry{Police station}
{/pəˈliːs ˌsteɪʃn/ (n)}
{The office or headquarters of a local police force.}
{The suspect was taken to the local police station for questioning.}
{Đây là một danh từ dùng để chỉ đồn cảnh sát.}

\vocabentry{Precedent}
{/ˈpresɪdənt/ (n)}
{An earlier event or action that is regarded as an example or guide to be considered in subsequent similar circumstances.}
{The judge's decision set a new precedent for how such cases are handled.}
{Đây là một danh từ dùng để chỉ tiền lệ pháp.}

\vocabentry{Prisoner}
{/ˈprɪznər/ (n)}
{A person legally committed to prison as a punishment or while awaiting trial.}
{The prisoner escaped from the high-security facility through a ventilation shaft.}
{Đây là một danh từ dùng để chỉ tù nhân.}

\vocabentry{Prosecute}
{/ˈprɑːsɪkjuːt/ (v)}
{Institute legal proceedings against (a person or organization).}
{The government has decided to prosecute the company for environmental pollution.}
{Động từ chỉ việc truy tố.}

\vocabentry{Punish}
{/ˈpʌnɪʃ/ (v)}
{Inflict a penalty or sanction on (someone) as retribution for an offense.}
{Parents should not use physical violence to punish their children.}
{Đây là một động từ dùng để chỉ việc trừng phạt.}

\vocabentry{Racketeering}
{/ˌrækɪˈtɪrɪŋ/ (n)}
{Dishonest and fraudulent business dealings.}
{The mob boss was finally convicted of racketeering and tax evasion.}
{Đây là một danh từ dùng để chỉ hành vi tống tiền hoặc làm ăn phi pháp có tổ chức.}

\vocabentry{Rape}
{/reɪp/ (n/v)}
{Unlawful sexual activity and usually sexual intercourse carried out forcibly.}
{The survivor of the rape showed incredible courage during the trial.}
{Đây là một danh từ/động từ dùng để chỉ tội hiếp dâm.}

\vocabentry{Rehabilitate}
{/ˌriːəˈbɪlɪteɪt/ (v)}
{Restore someone to health or normal life through training and therapy.}
{The program is designed to rehabilitate young offenders and help them find jobs.}
{Động từ chỉ việc cải tạo/phục hồi nhân phẩm.}

\vocabentry{Reprieve}
{/rɪˈpriːv/ (n/v)}
{Cancel or postpone the punishment of (someone, especially someone condemned to death).}
{The prisoner was granted a last-minute reprieve just hours before his execution.}
{Đây là một danh từ/động từ dùng để chỉ sự hoãn thi hành án.}

\vocabentry{Rob}
{/rɑːb/ (v)}
{Take property unlawfully from a person or place by force or threat of force.}
{They were planning to rob a bank but were caught by the police beforehand.}
{Đây là một động từ dùng để chỉ việc cướp.}

\vocabentry{Sanction}
{/ˈsæŋkʃn/ (n/v)}
{A threatened penalty for disobeying a law or rule.}
{The international community imposed sanctions on the country to stop the human rights abuses.}
{Đây là một danh từ/động từ dùng để chỉ lệnh trừng phạt hoặc việc xử phạt.}

\vocabentry{Serial killer}
{/ˌsɪriəl ˈkɪlər/ (n)}
{A person who commits a series of murders, often with no apparent motive and typically following a characteristic, predictable behavior pattern.}
{The police spent decades trying to capture the elusive serial killer.}
{Đây là một cụm danh từ dùng để chỉ kẻ sát nhân hàng loạt.}

\vocabentry{Shoplifter}
{/ˈʃɑːplɪftər/ (n)}
{A person who steals goods from a shop while pretending to be a customer.}
{Security guards caught a shoplifter trying to leave with a hidden bottle of perfume.}
{Danh từ chỉ kẻ trộm đồ trong cửa hàng.}

\vocabentry{Slander}
{/ˈslændər/ (n/v)}
{The action or crime of making a false spoken statement damaging to a person's reputation.}
{He sued his former employer for slander after they made false accusations about him.}
{Đây là một danh từ/động từ dùng để chỉ sự vu khống/phỉ báng bằng lời nói.}

\vocabentry{Smuggle}
{/ˈsmʌɡl/ (v)}
{Move (goods) illegally into or out of a country.}
{They were caught trying to smuggle rare birds out of the country in their luggage.}
{Động từ của smuggling, chỉ việc buôn lậu.}

\vocabentry{Solitary confinement}
{/ˌsɑːləteri kənˈfaɪnmənt/ (n)}
{The isolation of a prisoner in a separate cell as a punishment.}
{Many psychologists argue that long periods of solitary confinement can cause permanent mental damage.}
{Đây là một cụm danh từ dùng để chỉ hình phạt biệt giam.}

\vocabentry{Suit}
{/suːt/ (n)}
{A lawsuit.}
{She filed a civil suit against the driver who caused the accident.}
{Đây là một danh từ viết tắt dùng để chỉ vụ kiện.}

\vocabentry{Supreme Court}
{/suːˌpriːm ˈkɔːrt/ (n)}
{The highest judicial court in a country or state.}
{The Supreme Court's decision will have a major impact on civil rights in the country.}
{Đây là một cụm danh từ dùng để chỉ tòa án tối cao.}

\vocabentry{Theft}
{/θeft/ (n)}
{The action or crime of stealing.}
{Car theft has increased significantly in the suburban areas recently.}
{Lặp lại để nhấn mạnh nghĩa danh từ chỉ tội trộm cắp.}

\vocabentry{Tort}
{/tɔːrt/ (n)}
{A wrongful act or an infringement of a right (other than under contract) leading to civil legal liability.}
{Personal injury claims are a common type of tort in the legal system.}
{Đây là một danh từ dùng để chỉ trách nhiệm ngoài hợp đồng/lỗi dân sự.}

\vocabentry{Traffic}
{/ˈtræfɪk/ (v)}
{Deal or trade in something illegal. (Đây là một động từ dùng để chỉ hành vi buôn bán bất hợp pháp (người, ma túy...).).}
{The gang was involved in a global network that trafficked human organs.}
{}

\vocabentry{Unlawful}
{/ʌnˈlɔːfl/ (adj)}
{Not conforming to, permitted by, or recognized by law or rules.}
{The court ruled that the detention of the activists was unlawful and ordered their release.}
{Đây là một tính từ đồng nghĩa với illegal, chỉ sự trái luật.}

\vocabentry{Vice}
{/vaɪs/ (n)}
{Immoral or wicked behavior.}
{The police have launched a crackdown on vice in the city's red-light district.}
{Trong luật pháp thường dùng để chỉ các tệ nạn như mại dâm, cờ bạc.}

\vocabentry{Wrongful}
{/ˈrɔːŋfl/ (adj)}
{(Of an act) characterized by unfairness or injustice; contrary to law.}
{He received a large settlement for his wrongful dismissal from the company.}
{Đây là một tính từ dùng để chỉ tính sai trái hoặc bất công về mặt pháp lý.}

